
Grouped Query Attention (GQA) architectures share key-value (KV) heads across multiple query heads, introducing an aggregation decision that is implicit in all prior KV cache eviction methods: how to combine per-query-head importance scores into a single per-KV-head eviction score. We investigate whether max-pooling aggregation (GQA-GroupMax) over within-group query-head scores preserves minority-head signals that mean-pooling (GQA-GroupMean) suppresses due to a 1/G dilution factor. Two diagnostic experiments were conducted on LLaMA-3-8B, Mistral-7B-v0.1, and Phi-3-mini-4k-instruct using LongBench retrieval tasks. H-E1 measured within-group query-head attention cosine similarity to assess whether meaningful diversity exists within GQA groups. For the primary model LLaMA-3-8B, 40.6\% of layers exhibit median cosine similarity below 0.9, falling short of the pre-specified 50\% threshold; Mistral-7B (62.5\%) and Phi-3-mini (78.1\%) exceed the threshold. H-M1 measured Spearman rank correlation between GroupMax and GroupMean eviction score vectors on LLaMA-3-8B; the global median correlation was 0.962 (threshold for meaningful divergence: < 0.95), with 7 of 32 layers showing per-layer correlation below 0.95. Both gate criteria were not met on the primary model, constituting quantitative near-misses rather than decisive rejections. A key confound was identified: sequence truncation to 2048 tokens (required by memory constraints) reduces attention sparsity relative to full LongBench context lengths of 4,000--8,000 tokens, likely attenuating the GroupMax/GroupMean divergence. The aggregation function formalization and reusable measurement infrastructure are contributed regardless of gate outcome.

